\section{Limitations}

This is a standard-scale evidence-mapped review, not a comprehensive or systematic review. The 46-entry library passed bibliographic metadata audit, but only 25 local records include abstracts and 21 are title-only. Full-text inspection was selective, so a metadata PASS does not validate the context of every scientific claim. The manuscript consequently uses title-only entries for topic placement and labels abstract-reported associations at their actual provenance. It does not claim that the 46 records exhaust 2007--2026 LLZO work or that the small 2024--2026 subset represents current coverage.

The evidence matrix is intentionally condition-preserving rather than numerical. That choice prevents invented normalization but also means the review does not supply a performance meta-analysis. Missing fields are reported as unavailable in the checked source, not interpreted as experimental absence. The six gaps likewise diagnose what this collection cannot discriminate; full-text retrieval and a broader, reproducible search would be required before promoting them to field-wide absence claims.

The process map is a navigation asset derived from the adopted taxonomy. Its solid arrows denote process order and its dashed arrows denote boundary checks; neither encodes proven causality. The model-assisted box is more restricted still: its precursor sets are verified only as returned model output, its chemistry is unverified, its conditions are null, and the underlying reviewer issue I4 remains unresolved.

\section{Conclusion}

The central result is a disciplined way to read LLZO processing evidence. Powder synthesis, phase/composition control, thermal treatment/densification, and structure/transport form the primary chain, while composition, lithium inventory, environment, exposure, and additives cut across it. A separate evidence facet prevents those modifiers from being mistaken for stages or from disappearing behind route names.

The minimum comparability record converts that framework into a decision: results are directly comparable only when relevant composition, processing, structure, transport, and replication fields align; otherwise they remain parallel qualitative evidence or insufficiently specified. The six-gap map then points to three high-priority forms of discrimination: isolate starting lithium inventory from firing-boundary supply, separate density mediation from additive-related boundary chemistry, and harmonize density/transport records across specimens and laboratories.

These conclusions are calibrated to 46 audited records and selective abstract-level evidence. They do not rank a universally superior powder or densification route. Instead, Contributions A+B+C expose what was varied, what was measured, what traveled with the intervention, and what observation could overturn a working explanation. That is the level of synthesis the current evidence can support without turning missing fields, model outputs, or cross-paper incompatibility into chemical fact.
